\chapter{2020年浙江卷}

\section{单选题}

\begin{problem}
已知集合 \(P=\{x\mid 1<x<4\}\)，\(Q=\{x\mid 2<x<3\}\)，则 \(P\cap Q=\)（\quad）

\choices
    {\(\{x\mid 1<x\le 2\}\)}
    {\(\{x\mid 2<x<3\}\)}
    {\(\{x\mid 3\le x<4\}\)}
    {\(\{x\mid 1<x<4\}\)}

\end{problem}

\begin{answer}
B
\end{answer}

\begin{solution}
由交集的定义，
\(P\cap Q=\{x\mid 2<x<3\}\).
故选 B.
\end{solution}

\begin{problem}
已知 \(a\in\R\)，若 \(a-1+(a-2)\i\)（\(\i\) 为虚数单位）是实数，则 \(a=\)（\quad）

\choices
    {\(1\)}
    {\(-1\)}
    {\(2\)}
    {\(-2\)}

\end{problem}

\begin{answer}
C
\end{answer}

\begin{solution}
因为 \(a-1+(a-2)\i\) 是实数，所以虚部为 0，即
\(a-2=0\).
所以 \(a=2\). 故选 C.
\end{solution}

\begin{problem}
若实数 \(x\)，\(y\) 满足约束条件

\(\begin{cases}
x-3y+1\le0,\\
x+y-3\ge0,
\end{cases}\)

则 \(z=x+2y\) 的取值范围是（\quad）

\choices
    {\((-\infty,4]\)}
    {\([4,+\infty)\)}
    {\([5,+\infty)\)}
    {\((-\infty,+\infty)\)}

\end{problem}

\begin{answer}
B
\end{answer}

\begin{solution}
可行域由两条边界直线及其上方区域组成，顶点为 \((2,1)\).
把目标函数写成 \(z=x+2y\)，在可行域内向右上方平移时，\(z\) 递增，最小值在顶点 \((2,1)\) 处取得，
\(z_{\min}=2+2\times1=4\).
因此 \(z\) 的取值范围是 \([4,+\infty)\). 故选 B.
\end{solution}

\begin{problem}
函数 \(y=x\cos x+\sin x\) 在区间 \([-\pi,\pi]\) 的图象大致为（\quad）

\choices
    {\(\choicebitmap[width=0.15\paperwidth]{img/2020/zhejiang/q4_choice_a.png}\)}
    {\(\choicebitmap[width=0.15\paperwidth]{img/2020/zhejiang/q4_choice_b.png}\)}
    {\(\choicebitmap[width=0.15\paperwidth]{img/2020/zhejiang/q4_choice_c.png}\)}
    {\(\choicebitmap[width=0.15\paperwidth]{img/2020/zhejiang/q4_choice_d.png}\)}

\end{problem}

\begin{answer}
A
\end{answer}

\begin{solution}
设 \(f(x)=x\cos x+\sin x\)，则
\(f(-x)=-x\cos x-\sin x=-f(x)\)，
所以 \(f(x)\) 是奇函数，图象关于原点对称，排除 B，D.
又 \(f'(0)=2>0\)，且
\(f(\pi)=\pi\cos\pi+\sin\pi=-\pi<0\)，
故图象应与 A 一致. 故选 A.
\end{solution}

\begin{problem}
某几何体的三视图（单位：\(\mathrm{cm}\)）如图所示，则该几何体的体积（单位：\(\mathrm{cm}^3\)）是（\quad）

\begin{center}
\bitmapfigure[max width=0.46\linewidth]{img/2020/zhejiang/q5_fig1.png}
\end{center}

\choices
    {\(\dfrac{7}{3}\)}
    {\(\dfrac{14}{3}\)}
    {\(3\)}
    {\(6\)}

\end{problem}

\begin{answer}
A
\end{answer}

\begin{solution}
由三视图可知，该几何体由下部的直三棱柱和上部的三棱锥组成.
下部底面是斜边长为 \(2\) 的等腰直角三角形，面积为 \(1\)，高为 \(2\)，体积为 \(2\).
上部是底面积为 \(1\)，高为 \(1\) 的三棱锥，体积为 \(\dfrac13\).
所以几何体的体积为
\(2+\frac13=\frac73\).
故选 A.
\end{solution}

\begin{problem}
已知空间中不过同一点的三条直线 \(m\)，\(n\)，\(l\)，则“\(m\)，\(n\)，\(l\) 在同一平面”是“\(m\)，\(n\)，\(l\) 两两相交”的（\quad）

\choices
    {充分不必要条件}
    {必要不充分条件}
    {充分必要条件}
    {既不充分也不必要条件}

\end{problem}

\begin{answer}
B
\end{answer}

\begin{solution}
若 \(m\)，\(n\)，\(l\) 两两相交，则它们必在同一平面内，所以“三条直线在同一平面内”是“三条直线两两相交”的必要条件.
反过来，三条直线在同一平面内，并不一定两两相交，例如三条直线可以互相平行，所以不是充分条件.
故选 B.
\end{solution}

\begin{problem}
已知等差数列 \(\{a_n\}\) 的前 \(n\) 项和为 \(S_n\)，公差 \(d\neq0\)，且 \(\dfrac{a_1}{d}\le1\). 记 \(b_1=S_2\)，\(b_{n+1}=S_{2n+2}-S_{2n}\)（\(n\in\N^*\)），下列等式不可能成立的是（\quad）

\choices
    {\(2a_4=a_2+a_6\)}
    {\(2b_4=b_2+b_6\)}
    {\(a_4^2=a_2a_8\)}
    {\(b_4^2=b_2b_8\)}

\end{problem}

\begin{answer}
D
\end{answer}

\begin{solution}
设 \(r=\dfrac{a_1}{d}\)，则 \(a_n=d(r+n-1)\)，且 \(r\le1\).
由 \(S_n=\dfrac{nd}{2}\left(2r+n-1\right)\)，得
\(b_1=d(2r+1)\)，且 \(b_{n+1}=S_{2n+2}-S_{2n}=d(2r+4n+1)\).
等式 A 对任意等差数列都成立；等式 B 也恒成立.
等式 C 等价于
\((r+3)^2=(r+1)(r+7)\)，即 \(r=1\)，在题设范围内可以成立.
等式 D 等价于
\((2r+13)^2=(2r+5)(2r+29)\)，即 \(r=\dfrac32\)，与 \(r\le1\) 矛盾.
故选 D.
\end{solution}

\begin{problem}
已知点 \(O(0,0)\)，\(A(-2,0)\)，\(B(2,0)\). 设点 \(P\) 满足 \(|PA|-|PB|=2\)，且 \(P\) 为函数 \(y=3\sqrt{4-x^2}\) 图象上的点，则 \(|OP|=\)（\quad）

\choices
    {\(\dfrac{\sqrt{22}}{2}\)}
    {\(\dfrac{4\sqrt{10}}{5}\)}
    {\(\sqrt7\)}
    {\(\sqrt{10}\)}

\end{problem}

\begin{answer}
D
\end{answer}

\begin{solution}
由 \(|PA|-|PB|=2\) 可知，点 \(P\) 在以 \(A\)，\(B\) 为焦点的双曲线右支上，方程为
\(\frac{x^2}{1}-\frac{y^2}{3}=1\).
又 \(P\) 在 \(y=3\sqrt{4-x^2}\) 上，所以
\(y^2=9(4-x^2)\).
联立得
\(x^2-\frac{9(4-x^2)}{3}=1\)，
即
\(4x^2=13\).
于是
\(y^2=9\left(4-\frac{13}{4}\right)=\frac{27}{4}\)，
从而
\(|OP|^2=x^2+y^2=\frac{13}{4}+\frac{27}{4}=10\).
故 \(|OP|=\sqrt{10}\)，选 D.
\end{solution}

\begin{problem}
已知 \(a\)，\(b\in\R\) 且 \(ab\neq0\)，若 \((x-a)(x-b)(x-2a-b)\ge0\) 在 \(x\ge0\) 上恒成立，则（\quad）

\choices
    {\(a<0\)}
    {\(a>0\)}
    {\(b<0\)}
    {\(b>0\)}

\end{problem}

\begin{answer}
C
\end{answer}

\begin{solution}
设 \(f(x)=(x-a)(x-b)(x-2a-b)\).
若 \(b>0\)，则当 \(a>0\) 时，\(f(0)=-ab(2a+b)<0\)，矛盾；当 \(a<0\) 时，正数 \(b\) 是三次函数的一个正根，且它只能是单根，因为 \(a\ne b\)、\(a\ne0\)，\(2a+b\ne b\). 于是 \(f(x)\) 在 \(b\) 的两侧异号，与 \(x\ge0\) 时恒有 \(f(x)\ge0\) 矛盾.
所以 \(b<0\). 故选 C.
\end{solution}

\begin{problem}
设集合 \(S\)，\(T\)，\(S\subseteq\N^*\)，\(T\subseteq\N^*\)，\(S\)，\(T\) 中至少有两个元素，且 \(S\)，\(T\) 满足：

\begin{circlelist}
\item 对于任意 \(x\)，\(y\in S\)，若 \(x\ne y\)，都有 \(xy\in T\)；

\item 对于任意 \(x\)，\(y\in T\)，若 \(x<y\)，则 \(\dfrac{y}{x}\in S\).
\end{circlelist}

下列命题正确的是（\quad）

\choices
    {若 \(S\) 有 4 个元素，则 \(S\cup T\) 有 7 个元素}
    {若 \(S\) 有 4 个元素，则 \(S\cup T\) 有 6 个元素}
    {若 \(S\) 有 3 个元素，则 \(S\cup T\) 有 4 个元素}
    {若 \(S\) 有 3 个元素，则 \(S\cup T\) 有 5 个元素}

\end{problem}

\begin{answer}
A
\end{answer}

\begin{solution}
设 \(S=\{p_1,p_2,p_3,p_4\}\)，其中 \(p_1<p_2<p_3<p_4\). 由条件一和条件二，六个比值
\(\dfrac{p_2}{p_1}\)，\(\dfrac{p_3}{p_1}\)，\(\dfrac{p_4}{p_1}\)，\(\dfrac{p_3}{p_2}\)，\(\dfrac{p_4}{p_2}\)，\(\dfrac{p_4}{p_3}\)
都属于 \(S\).

若 \(p_1=1\)，则 \(1<p_3/p_2<p_3\)，所以 \(p_3/p_2=p_2\)，即 \(p_3=p_2^2\). 又
\(1<p_4/p_3<p_4/p_2<p_4\)，故 \(p_4/p_3=p_2\)，\(p_4/p_2=p_3\)，即 \(p_4=p_2^3\). 此时 \(T\) 中含有 \(p_2,p_2^2,p_2^3,p_2^4,p_2^5\)，由条件二得 \(p_2^4\in S\)，矛盾. 因而 \(p_1\ge2\).

此时 \(1<p_2/p_1<p_2\)，故 \(p_2/p_1=p_1\)，即 \(p_2=p_1^2\). 又
\(1<p_3/p_2<p_3/p_1<p_3\)，故 \(p_3/p_2=p_1\)，\(p_3/p_1=p_2\)，即 \(p_3=p_1^3\). 同理，由
\(1<p_4/p_3<p_4/p_2<p_4/p_1<p_4\)，得 \(p_4=p_1^4\).
令 \(p=p_1\)，则 \(S=\{p,p^2,p^3,p^4\}\)，且条件一给出
\(\{p^3,p^4,p^5,p^6,p^7\}\subseteq T\).
若 \(q\in T\) 且 \(q<p^3\)，则由 \(p^3/q\in S\) 可知 \(q\in\{1,p,p^2\}\)，但它们分别与 \(p^7\) 比较都会使商不属于 \(S\). 若 \(q>p^3\)，则 \(q/p^3\in S\)，故 \(q\in\{p^4,p^5,p^6,p^7\}\). 所以 \(T=\{p^3,p^4,p^5,p^6,p^7\}\)，从而 \(S\cup T\) 有 \(7\) 个元素，故选 A.
\end{solution}
\section{填空题}

\begin{problem}
已知数列 \(\{a_n\}\) 满足 \(a_n=\dfrac{n(n+1)}{2}\)，则 \(S_3=\fillinblank{}\).

\end{problem}

\begin{answer}
10
\end{answer}

\begin{solution}
\(a_1=1\)，\(a_2=3\)，\(a_3=6\)，
所以
\(S_3=1+3+6=10\).
\end{solution}

\begin{problem}
设 \((1+2x)^5=a_1+a_2x+a_3x^2+a_4x^3+a_5x^4+a_6x^5\)，则 \(a_5=\fillinblank{}\)；\(a_1+a_2+a_3=\fillinblank{}\).

\end{problem}

\begin{answer}
\(80\)，\(51\)
\end{answer}

\begin{solution}
\((1+2x)^5\) 的通项为
\(T_{r+1}=2^r\mathrm C_5^r x^r\).
令 \(r=4\)，得 \(a_5=2^4\mathrm C_5^4=80\)；又
\(a_1+a_2+a_3=2^0\mathrm C_5^0+2^1\mathrm C_5^1+2^2\mathrm C_5^2=1+10+40=51\).
\end{solution}

\begin{problem}
已知 \(\tan\theta=2\)，则 \(\cos2\theta=\fillinblank{}\)；\(\tan\left(\theta-\dfrac{\pi}{4}\right)=\fillinblank{}\).

\end{problem}

\begin{answer}
\(-\dfrac35\)，\(\dfrac13\)
\end{answer}

\begin{solution}
由 \(\tan\theta=2\)，得
\(\cos2\theta=\frac{1-\tan^2\theta}{1+\tan^2\theta}=\frac{1-4}{1+4}=-\frac35\).
又
\(\tan\left(\theta-\frac{\pi}{4}\right)=\frac{\tan\theta-1}{1+\tan\theta} =\frac{2-1}{1+2}=\frac13\).
\end{solution}

\begin{problem}
已知圆锥展开图的侧面积为 \(2\pi\)，且为半圆，则底面半径为\fillinblank{}.

\end{problem}

\begin{answer}
1
\end{answer}

\begin{solution}
设母线长为 \(l\). 因为展开图是半圆，所以其面积为
\(\frac12\pi l^2=2\pi\)，
得 \(l=2\).
半圆弧长就是底面圆的周长，故
\(2\pi r=\pi l=2\pi\)，
所以 \(r=1\).
\end{solution}

\begin{problem}
设直线 \(l:y=kx+b\)（\(k>0\)），圆 \(C_1:x^2+y^2=1\)，\(C_2:(x-4)^2+y^2=1\)，若直线 \(l\) 与 \(C_1\)，\(C_2\) 都相切，则 \(k=\fillinblank{}\)；\(b=\fillinblank{}\).

\end{problem}

\begin{answer}
\(\dfrac{\sqrt3}{3}\)，\(-\dfrac{2\sqrt3}{3}\)
\end{answer}

\begin{solution}
由切线到圆心的距离等于半径可得
\(\frac{|b|}{\sqrt{1+k^2}}=1\)，\(\frac{|4k+b|}{\sqrt{1+k^2}}=1\).
因 \(k>0\)，可取 \(b=-2k\). 代入上式得
\(4k^2=1+k^2\)，
故
\(k=\frac{\sqrt3}{3}\)，\(b=-2k=-\frac{2\sqrt3}{3}\).
\end{solution}

\begin{problem}
一个盒子里有 1 个红，1 个绿，2 个黄四个相同的球，每次拿一个，不放回，拿出红球即停，设拿出黄球的个数为 \(\xi\)，则 \(P(\xi=0)=\fillinblank{}\)；\(E(\xi)=\fillinblank{}\).

\end{problem}

\begin{answer}
\(\dfrac13\)，\(1\)
\end{answer}

\begin{solution}
在红球与两个黄球的相对顺序中，红球排在第 \(1\)、\(2\)、\(3\) 位的概率均为 \(\frac13\)，因此
\(P(\xi=0)=\frac13\)，\(P(\xi=1)=\frac13\)，\(P(\xi=2)=\frac13\).
所以
\(E(\xi)=0\times\frac13+1\times\frac13+2\times\frac13=1\).
\end{solution}

\begin{problem}
设 \(\bs e_1\)，\(\bs e_2\) 为单位向量，满足 \(|2\bs e_1-\bs e_2|\le\sqrt2\)，\(\bs a=\bs e_1+\bs e_2\)，\(\bs b=3\bs e_1+\bs e_2\)，设 \(\bs a\)，\(\bs b\) 的夹角为 \(\theta\)，则 \(\cos^2\theta\) 的最小值为\fillinblank{}.

\end{problem}

\begin{answer}
\(\dfrac{28}{29}\)
\end{answer}

\begin{solution}
设 \(\bs e_1\)，\(\bs e_2\) 的夹角为 \(\alpha\)，则
\(|2\bs e_1-\bs e_2|^2=4|\bs e_1|^2-4\bs e_1\cdot\bs e_2+|\bs e_2|^2=4-4\cos\alpha+1\le2\)，
从而
\(\cos\alpha\ge\frac34\).
又
\(\bs a\cdot\bs b=(\bs e_1+\bs e_2)\cdot(3\bs e_1+\bs e_2)=4+4\cos\alpha\)，
\(|\bs a|^2=|\bs e_1|^2+2\bs e_1\cdot\bs e_2+|\bs e_2|^2=2+2\cos\alpha\)，
\(|\bs b|^2=9|\bs e_1|^2+6\bs e_1\cdot\bs e_2+|\bs e_2|^2=10+6\cos\alpha\).
于是
\(\cos^2\theta=\frac{(\bs a\cdot\bs b)^2}{|\bs a|^2|\bs b|^2} =\frac{(4+4\cos\alpha)^2}{(2+2\cos\alpha)(10+6\cos\alpha)} =\frac{4+4\cos\alpha}{5+3\cos\alpha} =\frac43-\frac{\frac83}{5+3\cos\alpha}\).
所以当 \(\cos\alpha=\dfrac34\) 时，\(\cos^2\theta\) 取得最小值
\(\cos^2\theta_{\min}=\frac43-\frac{\frac83}{5+3\times\frac34}=\frac{28}{29}\).
\end{solution}

\section{解答题}

\begin{problem}
在锐角 \(\triangle ABC\) 中，角 \(A\)，\(B\)，\(C\) 的对边分别为 \(a\)，\(b\)，\(c\)，且 \(2b\sin A=\sqrt3\,a\).

\begin{enumerate}
\item 求角 \(B\)；

\item 求 \(\cos A+\cos B+\cos C\) 的取值范围.
\end{enumerate}
\end{problem}

\begin{solution}
\begin{enumerate}
\item 由正弦定理，\(\dfrac{a}{\sin A}=\dfrac{b}{\sin B}\)，所以 \(2b\sin A=\sqrt3\,a\) 化为 \(2\sin B\sin A=\sqrt3\,\sin A\). 因为 \(\sin A\ne0\)，故 \(\sin B=\dfrac{\sqrt3}{2}\). 又 \(B\) 为锐角，所以 \(B=\dfrac{\pi}{3}\).

\item 因为 \(B=\dfrac{\pi}{3}\)，所以 \(C=\pi-A-\dfrac{\pi}{3}=\dfrac{2\pi}{3}-A\). 于是
\(\cos A+\cos B+\cos C=\cos A+\dfrac12+\cos\left(\dfrac{2\pi}{3}-A\right)=\dfrac12+\cos\left(A-\dfrac{\pi}{3}\right)\).
又 \(\triangle ABC\) 为锐角三角形，所以 \(\dfrac{\pi}{6}<A<\dfrac{\pi}{2}\)，从而 \(-\dfrac{\pi}{6}<A-\dfrac{\pi}{3}<\dfrac{\pi}{6}\). 因此
\(\dfrac{\sqrt3}{2}<\cos\left(A-\dfrac{\pi}{3}\right)\le1\)，故
\(\cos A+\cos B+\cos C\in\left(\dfrac{1+\sqrt3}{2},\dfrac32\right]\).
\end{enumerate}
\end{solution}

\begin{problem}
如图，三棱台 \(DEF-ABC\) 中，面 \(ADFC\perp\) 面 \(ABC\)，\(\angle ACB=\angle ACD=45^\circ\)，\(DC=2BC\).

\begin{enumerate}
\item 证明：\(EF\perp DB\)；

\item 求 \(DF\) 与面 \(DBC\) 所成角的正弦值.
\end{enumerate}

\begin{center}
\bitmapfigure[max width=0.36\linewidth]{img/2020/zhejiang/q19_fig1.png}
\end{center}
\end{problem}

\begin{solution}
\begin{enumerate}
\item 证明：作 \(DH\perp AC\)，交 \(AC\) 于点 \(H\).
因为面 \(ADFC\perp\) 面 \(ABC\)，\(DH\subset\) 面 \(ADFC\)，所以 \(DH\perp BC\).
在 \(\mathrm{Rt}\triangle DHC\) 中，
\(CH=CD\cos45^\circ=\frac{\sqrt2}{2}CD\).
因为 \(DC=2BC\)，所以
\(CH=\frac{\sqrt2}{2}CD=\frac{\sqrt2}{2}\cdot2BC=\sqrt2\cdot BC\).
所以
\(\frac{BC}{CH}=\frac{\sqrt2}{2}\)，
即 \(\triangle BHC\) 是直角三角形，且 \(\angle HBC=90^\circ\).
所以 \(HB\perp BC\)，所以 \(BC\perp\) 面 \(DHB\)，因为 \(BD\subset\) 面 \(DHB\)，所以 \(BC\perp BD\).
因为在三棱台 \(DEF-ABC\) 中，\(EF//BC\)，所以 \(EF\perp DB\).

\item 设 \(BC=1\)，则 \(BH=1\)，\(HC=\sqrt2\).
在 \(\mathrm{Rt}\triangle DHC\) 中，\(DH=\sqrt2\)，\(DC=2\).
在 \(\mathrm{Rt}\triangle DHB\) 中，
\(DB=\sqrt{DH^2+HB^2}=\sqrt{2+1}=\sqrt3\).
作 \(HG\perp BD\) 于 \(G\)，因为 \(BC\perp HG\)，所以 \(HG\perp\) 面 \(BCD\)，因为 \(GC\subset\) 面 \(BCD\)，所以 \(HG\perp GC\)，所以 \(\triangle HGC\) 是直角三角形，且 \(\angle HGC=90^\circ\).
设 \(DF\) 与面 \(DBC\) 所成角为 \(\theta\)，则 \(\theta\) 即为 \(CH\) 与面 \(DBC\) 的夹角，且
\(\sin\theta=\sin\angle HCG=\frac{HG}{HC}=\frac{HG}{\sqrt2}\).
因为在 \(\mathrm{Rt}\triangle DHB\) 中，\(DH\cdot HB=BD\cdot HG\)，所以
\(HG=\frac{DH\cdot HB}{BD}=\frac{\sqrt2\cdot1}{\sqrt3}=\frac{\sqrt6}{3}\).
所以
\(\sin\theta=\frac{HG}{\sqrt2}=\frac{\sqrt6/3}{\sqrt2}=\frac{\sqrt3}{3}\).
故所求正弦值为 \(\dfrac{\sqrt3}{3}\).
\end{enumerate}
\end{solution}

\begin{problem}
已知数列 \(\{a_n\}\)，\(\{b_n\}\)，\(\{c_n\}\) 中，\(a_1=b_1=c_1=1\)，\(c_{n+1}=a_{n+1}-a_n\)，\(c_{n+1}=\dfrac{b_n}{b_{n+2}}\,c_n\)（\(n\in\N^*\)）.

\begin{enumerate}
\item 若数列 \(\{b_n\}\) 为等比数列，且公比 \(q>0\)，且 \(b_1+b_2=6b_3\)，求 \(q\) 与 \(a_n\) 的通项公式；

\item 若数列 \(\{b_n\}\) 为等差数列，且公差 \(d>0\)，证明：\(c_1+c_2+\cdots+c_n<1+\dfrac1d\)，\(n\in\N^*\).
\end{enumerate}
\end{problem}

\begin{solution}
\begin{enumerate}
\item 因为 \(\{b_n\}\) 为等比数列，所以 \(b_2=q\)，\(b_3=q^2\). 由 \(b_1+b_2=6b_3\) 得 \(1+q=6q^2\)，所以 \(q=\dfrac12\)（另一个根为 \(-\dfrac13\)，舍去）.
又 \(c_{n+1}=\dfrac{b_n}{b_{n+2}}c_n=\dfrac{1}{q^2}c_n=4c_n\)，且 \(c_1=1\)，故 \(c_n=4^{n-1}\). 于是 \(a_{n+1}-a_n=c_{n+1}=4^n\)，累加得
\(a_n=1+4+4^2+\cdots+4^{n-1}=\dfrac{4^n-1}{3}\).

\item 若 \(\{b_n\}\) 为等差数列且公差 \(d>0\)，则 \(b_n=1+(n-1)d\). 由 \(c_{n+1}=\dfrac{b_n}{b_{n+2}}c_n\) 可得 \(b_n b_{n+1}c_n\) 为常数. 由 \(b_1=c_1=1\) 得该常数为 \(1+d\)，所以
\(c_n=\dfrac{1+d}{b_n b_{n+1}}=\left(1+\dfrac1d\right)\left(\dfrac1{b_n}-\dfrac1{b_{n+1}}\right)\).
因此
\(c_1+c_2+\cdots+c_n=\left(1+\dfrac1d\right)\left(1-\dfrac1{b_{n+1}}\right)<1+\dfrac1d\).
\end{enumerate}
\end{solution}

\begin{problem}
如图，已知椭圆 \(C_1:\dfrac{x^2}{2}+y^2=1\)，抛物线 \(C_2:y^2=2px\)（\(p>0\)），点 \(A\) 是椭圆 \(C_1\) 与抛物线 \(C_2\) 的交点，过点 \(A\) 的直线 \(l\) 交椭圆 \(C_1\) 于点 \(B\)，交抛物线 \(C_2\) 于 \(M\)（\(B\)，\(M\) 不同于 \(A\)）.

\begin{enumerate}
\item 若 \(p=\dfrac{1}{16}\)，求抛物线 \(C_2\) 的焦点坐标；

\item 若存在不过原点的直线 \(l\) 使 \(M\) 为线段 \(AB\) 的中点，求 \(p\) 的最大值.
\end{enumerate}

\begin{center}
\bitmapfigure[max width=0.38\linewidth]{img/2020/zhejiang/q21_fig1.png}
\end{center}
\end{problem}

\begin{solution}
\begin{enumerate}
\item \(p=\dfrac1{16}\)，则 \(\dfrac p2=\dfrac1{32}\)，所以抛物线 \(C_2\) 的焦点坐标为 \(\left(\dfrac{1}{32},0\right)\).

\item 直线 \(l\) 与 \(x\) 轴垂直时，此时点 \(M\) 与点 \(A\) 或点 \(B\) 重合，不满足题意，
设直线 \(l\) 的方程为 \(y=kx+t\)，\(A(x_1,y_1)\)，\(B(x_2,y_2)\)，\(M(x_0,y_0)\)，
\examdisplaycases{}{
\dfrac{x^2}{2}+y^2=1,\\
y=kx+t
}
消去 \(y\) 可得
\((2k^2+1)x^2+4ktx+2t^2-2=0\)，
所以
\(\Delta=16k^2t^2-4(2k^2+1)(2t^2-2)\ge0\)，
即
\(t^2<1+2k^2\).
又
\(x_1+x_2=-\frac{4kt}{1+2k^2}\)，
所以
\(x_0=\frac12(x_1+x_2)=-\frac{2kt}{1+2k^2}\)，
\(y_0=kx_0+t=\frac{t}{1+2k^2}\)，
所以
\(M\left(-\frac{2kt}{1+2k^2},\frac{t}{1+2k^2}\right)\).
因为点 \(M\) 在抛物线 \(C_2\) 上，所以 \(y^2=2px\)，
\(p=\frac{y_0^2}{2x_0} =\frac{\dfrac{t^2}{(1+2k^2)^2}}{2\cdot\dfrac{-2kt}{1+2k^2}} =\frac{t}{-4k(1+2k^2)}\).
\examdisplaycases{}{
y^2=2px,\\
y=kx+t
}
解得
\(x_1=\frac{t(1+2k^2)}{-2k^3}\)，\(y_1=\frac{t}{-2k^2}\).
代入椭圆方程可得
\(\frac{t^2(1+2k^2)^2}{8k^6}+\frac{t^2}{4k^4}=1\)，
解得
\(t^2=\frac{8k^6}{(1+2k^2)^2+2k^2}\).
所以
\(p^2=\frac{t^2}{16k^2(1+2k^2)^2} =\frac{8k^6}{16k^2(1+2k^2)^2[(1+2k^2)^2+2k^2]} =\frac{k^4}{2(1+2k^2)^2[(1+2k^2)^2+2k^2]}\).
由基本不等式可得
\(p^2\le\frac{k^4}{2(2\sqrt2k)^2[(2\sqrt2k)^2+2k^2]}=\frac1{160}\).
所以
\(p\le\frac{\sqrt{10}}{40}\).
当且仅当 \(1=2k^2\)，即 \(k^2=\dfrac12\)，\(t^2=\dfrac15\) 时等号成立，
故 \(p\) 的最大值为 \(\dfrac{\sqrt{10}}{40}\).
\end{enumerate}
\end{solution}

\begin{problem}
已知 \(1<a\le2\)，函数 \(f(x)=\e^x-x-a\)，其中 \(\e=2.71828\cdots\) 为自然对数的底数.

\begin{enumerate}
\item 证明：函数 \(y=f(x)\) 在 \((0,+\infty)\) 上有唯一零点；

\item 记 \(x_0\) 为函数 \(y=f(x)\) 在 \((0,+\infty)\) 上的零点，证明：

\begin{enumerate}
\item \(\sqrt{a-1}\le x_0\le\sqrt{2(a-1)}\)；
\item \(x_0f(\e^{x_0})\ge(\e-1)(a-1)a\).
\end{enumerate}
\end{enumerate}
\end{problem}

\begin{solution}
\begin{enumerate}
\item 当 \(x>0\) 时，\(f'(x)=\e^x-1>0\)，所以 \(f(x)\) 在 \((0,+\infty)\) 上单调递增. 又 \(f(0)=1-a<0\)，且 \(f(2)=\e^2-2-a\ge\e^2-4>0\)，所以函数 \(y=f(x)\) 在 \((0,+\infty)\) 上有唯一零点.

\item
\begin{enumerate}
\item 由 \(f(x_0)=0\) 得 \(a-1=\e^{x_0}-x_0-1\). 因为 \(\e^x\ge1+x+\dfrac{x^2}{2}\)，所以 \(a-1\ge\dfrac{x_0^2}{2}\)，从而 \(x_0\le\sqrt{2(a-1)}\).

设 \(t\) 为 \(a=2\) 时方程 \(\e^x-x-2=0\) 的正根，则 \(x_0\le t\)，且 \(f_2(1)=\e-3<0\)，故 \(t>1\). 记 \(h(x)=\e^x-1-x-x^2\). 在 \(0\le x\le t\) 上，\(h''(x)=\e^x-2\)，所以 \(h'(x)\) 先减后增；又 \(h'(0)=0\)，\(h'(t)=\e^t-1-2t=1-t<0\)，故 \(h'(x)\le0\)，从而 \(h(x)\le h(0)=0\). 因此 \(\e^{x_0}-x_0-1\le x_0^2\)，即 \(a-1\le x_0^2\)，所以 \(x_0\ge\sqrt{a-1}\). 综上，\(\sqrt{a-1}\le x_0\le\sqrt{2(a-1)}\).

\item 令 \(s=\sqrt{a-1}\). 由上一问，\(x_0\ge s\)，且 \(f\) 单调递增、\(f(\e^{x_0})>0\)，所以
\(x_0f(\e^{x_0})=x_0f(x_0+a)\ge s f(s+a)=s\left(\e^{s+a}-s-2a\right)\).
由 \(\e^x\ge1+x+\dfrac{x^2}{2}\)，只需证明
\(1+\dfrac12(s+a)^2-a\ge(\e-1)as\).
由于 \(s^2=a-1\)，上式左边减去右边等于
\(\dfrac12\left(a^2-s^2-2(\e-2)as\right)\).
而
\(\dfrac{a}{s}=s+\dfrac1s\ge2\)，所以
\(\dfrac{a}{s}-\dfrac{s}{a}\ge2-\dfrac12=\dfrac32>2(\e-2)\).
故 \(a^2-s^2-2(\e-2)as\ge0\)，从而
\(x_0f(\e^{x_0})\ge(\e-1)(a-1)a\).
\end{enumerate}
\end{enumerate}
\end{solution}
